\chapter{Variational Monte Carlo}
\label{chap:vmc}
%=============================================================================

Chapter~\ref{chap:statelements} assembled the tools -- probability
distributions, the central limit theorem, the correlation time, Markov chains
and the Metropolis algorithm -- and ended with a first variational Monte Carlo
calculation, run brute force.  This chapter is the method itself, worked
through properly and taken to the point where it is usable: the variational
principle it rests on, how a trial wave function is constructed and why it has
the form it does, the local energy and how to differentiate it, the sampling
algorithm, and finally importance sampling, which is what turns a working
program into an efficient one.

The example throughout is the system of Chapter~\ref{chap:applications}, two
electrons in a two-dimensional parabolic trap.  It is the right test case for
three reasons.  It is small enough that everything can be differentiated by
hand; it has an analytic answer, $E=3$ at $\hbar\omega=1$, so we always know
what we are aiming at; and it has already been solved by Hartree-Fock,
perturbation theory and coupled cluster in Chapter~\ref{chap:applications}, so
the Monte Carlo result can be placed on the same scale as those.

The programs are in \texttt{vmc.py}.  Everything the sampling needs -- the
trial function, the local energy and the quantum force -- is available in
closed form, and each closed form is checked against finite differences before
it is used, for a reason that will become clear in
Section~\ref{sec:vmcimportance}.

%----------------------------------------------------------------------
\section{The variational principle}
\label{sec:vmcprinciple}

\index{variational Monte Carlo}\index{variational principle}
Everything rests on one inequality.  For any normalisable trial state
$\Psi_T$,
\begin{equation}
  \boxed{\;
  E[\Psi_T]
   = \frac{\bracket{\Psi_T}{\hat{H}}{\Psi_T}}{\braket{\Psi_T}{\Psi_T}}
   \;\ge\; E_0 \;}
  \label{eq:13-variational}
\end{equation}
with equality only if $\Psi_T$ is the ground state.  The proof is the one from
Section~\ref{sec:variational}: expand $\Psi_T$ in the exact eigenstates and
observe that every excited state can only raise the average.

The strategy is then obvious.  Choose a family of trial functions
$\Psi_T(\bm{R};\bm{\theta})$ depending on a few parameters $\bm{\theta}$,
evaluate $E[\Psi_T]$, and minimise over $\bm{\theta}$.  What makes this hard
is the evaluation.  For $N$ particles in $d$ dimensions the integral
\begin{equation}
  E[\Psi_T]
   = \frac{\int d\bm{R}\,\Psi_T^{*}(\bm{R})\hat{H}\Psi_T(\bm{R})}
          {\int d\bm{R}\,|\Psi_T(\bm{R})|^2}
  \label{eq:13-energyintegral}
\end{equation}
is over $Nd$ dimensions.  For our two electrons in two dimensions that is
four, which a grid could still manage; for the six-electron dot it is twelve,
and for anything realistic it is hopeless.  Monte Carlo integration converges
as $1/\sqrt{M}$ regardless of dimension, which is the whole reason for the
method.

\paragraph{What the method can and cannot do.}
It is worth being clear about the trade before starting.

\begin{svgraybox}
\begin{itemize}
\item Any expectation value can be computed, for any Hamiltonian, in any
  number of dimensions, provided we can write down and evaluate a trial wave
  function.  There is no basis set and therefore no basis-set error.
\item The energy is a rigorous upper bound, and the bound improves
  monotonically as the trial function improves.
\item But the answer is only as good as the trial function.  Unlike coupled
  cluster or configuration interaction there is no systematic hierarchy to
  climb; improving $\Psi_T$ is a matter of physical insight.
\item And the error is statistical, so an extra digit costs a hundred times
  the computer time.
\end{itemize}
\end{svgraybox}

%----------------------------------------------------------------------
\section{The local energy}
\label{sec:vmclocalenergy}

\index{local energy}
The step that makes Eq.~(\ref{eq:13-energyintegral}) computable was taken in
Section~\ref{sec:mcintegration}.  Insert $\Psi_T/\Psi_T$ and read the result
as an average over a probability density,
\begin{equation}
  E[\Psi_T]
   = \int d\bm{R}\;
     \underbrace{\frac{|\Psi_T(\bm{R})|^2}
       {\int d\bm{R}'|\Psi_T(\bm{R}')|^2}}_{\textstyle P(\bm{R})}\;
     \underbrace{\frac{\hat{H}\Psi_T(\bm{R})}{\Psi_T(\bm{R})}}_{
       \textstyle E_L(\bm{R})},
  \label{eq:13-localenergyaverage}
\end{equation}
where $P$ is the quantum-mechanical probability density and
\begin{equation}
  \boxed{\;E_L(\bm{R}) = \frac{1}{\Psi_T(\bm{R})}\hat{H}\Psi_T(\bm{R})\;}
  \label{eq:13-localenergy}
\end{equation}
is the \emph{local energy}.  The energy is then simply
\begin{equation}
  E[\Psi_T] \approx \frac{1}{M}\sum_{i=1}^{M}E_L(\bm{R}_i),
  \qquad \bm{R}_i \sim P ,
  \label{eq:13-energyestimator}
\end{equation}
an average of a function over samples drawn from $P$.  Two observations make
this more than a rewriting.

\paragraph{The zero-variance principle.}
If $\Psi_T$ is an exact eigenstate, $\hat{H}\Psi_T=E\Psi_T$, then
$E_L(\bm{R})=E$ everywhere.  The average is $E$ exactly, from a single
sample, and the variance is zero.  For an approximate $\Psi_T$ the variance
\begin{equation}
  \sigma^2 = \mean{E_L^2} - \mean{E_L}^2
  \label{eq:13-variance}
\end{equation}
is a direct and absolute measure of how far the trial function is from an
eigenstate -- absolute because we know the target value, zero.  This is why
$\sigma^2$ is quoted alongside the energy in everything that follows, and why
minimising the variance is a legitimate alternative to minimising the energy.

\paragraph{The normalisation never appears.}
$P$ contains the normalisation integral we cannot compute.  But the Metropolis
algorithm of Section~\ref{sec:metropolis} needs only ratios
$P(\bm{R}')/P(\bm{R}) = |\Psi_T(\bm{R}')|^2/|\Psi_T(\bm{R})|^2$, in which the
normalisation cancels.  That is the second reason the whole scheme works, and
it is why a Markov chain is used instead of direct sampling.

%----------------------------------------------------------------------
\section{Constructing a trial wave function}
\label{sec:vmctrial}

\index{trial wave function}\index{cusp condition}
A trial wave function is usually built as a product of two pieces,
\begin{equation}
  \Psi_T(\bm{R}) = \Phi(\bm{R})\,J(\bm{R}),
  \label{eq:13-productansatz}
\end{equation}
a mean-field part $\Phi$ -- a Slater determinant of single-particle orbitals,
which enforces antisymmetry -- and a symmetric \emph{correlation} or
\emph{Jastrow} factor $J$ depending on the interparticle distances.  The
determinant carries the shell structure; the Jastrow factor carries what a
determinant cannot represent, namely that two electrons avoid each other.

The form of $J$ is not guesswork.  It is fixed at short distance by a
condition on the local energy.  As two electrons approach, the Coulomb term
$1/r_{12}$ in $\hat{H}$ diverges, and if the local energy is to remain finite
the kinetic term must produce a compensating divergence.  Isolating the
relative motion and keeping only the terms that survive as $r_{12}\to0$, the
requirement $E_L$ finite gives
\begin{equation}
  \left(\frac{2}{r_{12}}\frac{d}{dr_{12}}
   + \frac{2}{r_{12}}\right){\cal R}(r_{12}) = 0
  \quad\Longrightarrow\quad
  \frac{1}{{\cal R}}\frac{d{\cal R}}{dr_{12}}
   = \frac{1}{2(l+1)}\bigg|_{r_{12}\to0},
  \label{eq:13-cuspderivation}
\end{equation}
the \emph{cusp condition}.  For $l=0$ and two electrons of opposite spin this
integrates to
\begin{equation}
  {\cal R}(r_{12}) \propto \exp\left(\frac{r_{12}}{2}\right)
  \qquad\text{(antiparallel spins)},
  \label{eq:13-cusp3d}
\end{equation}
in three dimensions, and to $\exp(r_{12}/4)$ for parallel spins, where the
Pauli principle already keeps the electrons apart and less correlation is
needed.  In two dimensions the corresponding coefficients are $1$ and
$1/3$.  For more than two electrons the condition holds pair by pair.

The exponential in Eq.~(\ref{eq:13-cusp3d}) is right at short distance and
badly wrong at long distance, where it grows without bound.  The standard
repair is the \emph{Pad\'e-Jastrow} form, which interpolates between the cusp
behaviour and a constant,
\begin{equation}
  J = \exp\left(\sum_{i<j}\frac{a_{ij}r_{ij}}{1+\beta r_{ij}}\right),
  \label{eq:13-padejastrow}
\end{equation}
with $a_{ij}$ fixed by the cusp condition and $\beta$ variational.  As
$r_{ij}\to0$ the exponent is $a_{ij}r_{ij}$ and the cusp is reproduced; as
$r_{ij}\to\infty$ it saturates at $a_{ij}/\beta$.

\begin{notebox}
\textbf{Why this matters later.} The construction above is physics put in by
hand: we knew to expect a cusp, we knew the exponential was wrong at long
range, and we chose a functional form with the right limits.  It works well
for two electrons and less well as systems grow, because the correlations we
have to guess become more complicated.  That difficulty is precisely the
motivation for the machine-learning wave functions of the later chapters,
where the functional form is not guessed but learned.
\end{notebox}

%----------------------------------------------------------------------
\section{The two-electron quantum dot}
\label{sec:vmcdot}

\index{quantum dot!variational Monte Carlo}
The system is the one of Chapter~\ref{chap:applications}: two electrons
confined to a plane by a parabolic trap, interacting through the bare Coulomb
repulsion.  In natural units,
\begin{equation}
  \hat{H} = \sum_{i=1}^{2}
    \left(-\tfrac12\nabla_i^2 + \tfrac12\omega^2 r_i^2\right)
    + \frac{1}{r_{12}} .
  \label{eq:13-hamiltonian}
\end{equation}

\paragraph{Separating the motion.}
Introducing the centre of mass $\bm{R}=(\bm{r}_1+\bm{r}_2)/2$ and the relative
coordinate $\bm{r}=\bm{r}_1-\bm{r}_2$, the Hamiltonian separates.  The
centre-of-mass part is a plain harmonic oscillator and contributes a known
constant; the relative part is
\begin{equation}
  \hat{H}_r = -\nabla_r^2 + \tfrac14\omega^2 r^2 + \frac{1}{r},
  \label{eq:13-relative}
\end{equation}
a one-body problem in the relative coordinate.  Taut showed that this has
closed-form solutions for a discrete set of frequencies, and at
$\hbar\omega=1$ the total ground-state energy is exactly $3$.  That number is
the target for the rest of the chapter.

\paragraph{The trial function.}
Combining the mean-field piece -- for two electrons in a spin singlet, simply
the product of two lowest oscillator orbitals -- with the Pad\'e-Jastrow
factor gives
\begin{equation}
  \boxed{\;
  \Psi_T(\bm{r}_1,\bm{r}_2;\alpha,\beta)
   = \exp\left(-\frac{\alpha\omega}{2}\left(r_1^2+r_2^2\right)\right)
     \exp\left(\frac{r_{12}}{1+\beta r_{12}}\right) \;}
  \label{eq:13-trial}
\end{equation}
Two variational parameters.  The first, $\alpha$, scales the width of the
orbital: $\alpha=1$ is the bare oscillator ground state, and letting it vary
allows the orbital to relax in response to the repulsion.  The second,
$\beta$, sets the range over which the correlation factor acts.  The numerator
of the Jastrow exponent is \emph{not} a parameter -- it is fixed at $1$ by the
cusp condition.

\paragraph{The local energy.}
Differentiating Eq.~(\ref{eq:13-trial}) twice and dividing by $\Psi_T$ gives,
with $d = 1/(1+\beta r_{12})$,
\begin{equation}
  \boxed{
  \begin{aligned}
    E_L =\;& \frac{1}{2}\omega^2\left(1-\alpha^2\right)
        \left(r_1^2+r_2^2\right) + 2\alpha\omega + \frac{1}{r_{12}}\\
      &+ d^2\left(\alpha\omega r_{12} - d^2 + 2\beta d
        - \frac{1}{r_{12}}\right).
  \end{aligned}}
  \label{eq:13-localenergydot}
\end{equation}
The first line is the result for the pure Gaussian, already met as
Eq.~(\ref{eq:12-alphaenergy})'s integrand in
Chapter~\ref{chap:statelements}; the second is everything the Jastrow factor
contributes.  Note the last term.  It carries $-d^2/r_{12}$, which at short
distance cancels the $+1/r_{12}$ on the first line, since $d\to1$ as
$r_{12}\to0$.  That cancellation \emph{is} the cusp condition, and
Table~\ref{tab:13-cusp} shows it happening.

\begin{table}[h]
\centering
\renewcommand{\arraystretch}{1.25}
\setlength{\tabcolsep}{16pt}
\begin{tabular}{rll}
\toprule
$r_{12}$ & $E_L$ with Jastrow & $E_L$ without\\
\midrule
$1.0000$ & $3.031237$ & $\phantom{0000}3.00$\\
$0.1000$ & $2.703286$ & $\phantom{000}12.00$\\
$0.0100$ & $2.611464$ & $\phantom{00}102.00$\\
$0.0010$ & $2.601159$ & $\phantom{0}1002.00$\\
$0.0001$ & $2.600116$ & $10002.00$\\
\bottomrule
\end{tabular}
\caption{The local energy for two electrons brought together along a line,
$\alpha=1$, $\beta=0.4$, $\hbar\omega=1$.  With the correlation factor the
local energy converges; without it the Coulomb singularity is unopposed and
$E_L$ diverges as $1/r_{12}$.  Computed by \texttt{vmc.py}.}
\label{tab:13-cusp}
\end{table}

This is worth pausing on, because it explains the numbers in
Chapter~\ref{chap:statelements}.  A trial function without the cusp has a
local energy with an unbounded tail; rare close encounters then dominate the
variance, the variance estimate itself becomes unstable, and error bars stop
meaning much.  With the cusp built in, the variance for this system falls by
three orders of magnitude.

%----------------------------------------------------------------------
\section{Brute-force Metropolis sampling}
\label{sec:vmcbrute}

\index{Metropolis algorithm!in variational Monte Carlo}
The algorithm is the Metropolis of Section~\ref{sec:metropolis} applied to
$P=|\Psi_T|^2$.

\begin{svgraybox}
\begin{enumerate}
\item Place the particles at random and evaluate $\Psi_T$.
\item Propose a move of \emph{one} particle,
  $\bm{r}_i' = \bm{r}_i + \Delta\cdot(\bm{u}-\tfrac12)$, with $\bm{u}$
  uniform and $\Delta$ the step length.
\item Accept with probability
  $\min\left(1,\;|\Psi_T(\bm{R}')|^2/|\Psi_T(\bm{R})|^2\right)$; if the move
  is rejected, keep the old configuration.
\item After every particle has been offered a move, accumulate $E_L(\bm{R})$
  and $E_L^2(\bm{R})$.
\item Discard an initial burn-in, then average.
\end{enumerate}
\end{svgraybox}

\noindent
Two details are easy to get wrong.  Moving one particle at a time rather than
all of them keeps the acceptance rate up for a given step length.  And a
rejected move must still be counted -- the old configuration is measured
again.  Dropping rejected steps biases the average, and is one of the most
common bugs in a first implementation.

The step length $\Delta$ is the one free parameter, and
Table~\ref{tab:13-brute} shows the usual trade-off: too small and the walker
crawls, too large and everything is rejected, with the correlation time large
at both ends.

\begin{table}[h]
\centering
\renewcommand{\arraystretch}{1.25}
\setlength{\tabcolsep}{11pt}
\begin{tabular}{rrlll}
\toprule
$\Delta$ & acceptance & energy & error & $\tau$\\
\midrule
$0.5$ & $90.1\%$ & $2.999609$ & $0.001249$ & $24.28$\\
$1.0$ & $80.2\%$ & $2.999422$ & $0.000643$ & $\phantom{0}6.17$\\
$2.0$ & $62.2\%$ & $3.000750$ & $0.000418$ & $\phantom{0}2.80$\\
$4.0$ & $35.3\%$ & $2.999901$ & $0.000376$ & $\phantom{0}2.24$\\
$8.0$ & $11.8\%$ & $3.000125$ & $0.000701$ & $\phantom{0}7.77$\\
\bottomrule
\end{tabular}
\caption{Brute-force Metropolis for the two-electron dot at $\alpha=0.98$,
$\beta=0.40$, $\hbar\omega=1$, $30\,000$ cycles.  Errors are corrected for the
autocorrelation time by Eq.~(\ref{eq:12-tau}).  Computed by \texttt{vmc.py}.}
\label{tab:13-brute}
\end{table}

%----------------------------------------------------------------------
\section{Importance sampling}
\label{sec:vmcimportance}

\index{importance sampling!in variational Monte Carlo}
The uniform proposal of the previous section is wasteful.  It offers moves in
every direction with equal probability, including many into regions where
$|\Psi_T|^2$ is negligible, and those are rejected.  A better proposal would
know something about the wave function and bias the move towards where it is
large.  The machinery for that is the Fokker-Planck and Langevin equations of
Section~\ref{sec:fokkerplanck}.

\paragraph{Drift from the Fokker-Planck equation.}
A diffusion process with a drift obeys
\begin{equation}
  \frac{\partial P}{\partial t}
   = \sum_i D\frac{\partial}{\partial \bm{x}_i}
     \left(\frac{\partial}{\partial \bm{x}_i} - \bm{F}_i\right)P(\bm{x},t),
  \label{eq:13-fokkerplanck}
\end{equation}
with $D$ the diffusion coefficient and $\bm{F}$ the drift.  We want the
stationary solution to be our target $P=|\Psi_T|^2$.  Setting the left-hand
side to zero and requiring every term of the sum to vanish gives
\begin{equation}
  \frac{\partial^2 P}{\partial \bm{x}_i^2}
   = P\frac{\partial \bm{F}_i}{\partial \bm{x}_i}
   + \bm{F}_i\frac{\partial P}{\partial \bm{x}_i} .
  \label{eq:13-stationary}
\end{equation}
Trying a drift of the form $\bm{F}=g(\bm{x})\,\partial P/\partial\bm{x}$ and
demanding that the first- and second-derivative terms cancel forces
$g=1/P$, so that $\bm{F}=\nabla P/P$.  With $P=|\Psi_T|^2$ this is
\begin{equation}
  \boxed{\;\bm{F} = \frac{2}{\Psi_T}\nabla\Psi_T = 2\nabla\ln\Psi_T\;}
  \label{eq:13-quantumforce}
\end{equation}
the \emph{quantum force}.  It points towards regions where the trial function
is large, and it is what replaces the aimless uniform proposal.

\paragraph{Proposing a move.}
The Langevin equation corresponding to Eq.~(\ref{eq:13-fokkerplanck}) is
\begin{equation}
  \frac{\partial \bm{x}(t)}{\partial t} = D\bm{F}(\bm{x}(t)) + \bm{\eta},
  \label{eq:13-langevin}
\end{equation}
with $\bm{\eta}$ a Gaussian white noise.  Integrating it over a short time
step $\Delta t$ by Euler's method gives the proposal
\begin{equation}
  \boxed{\;
  \bm{y} = \bm{x} + D\bm{F}(\bm{x})\Delta t
    + \bm{\xi}\sqrt{\Delta t} \;}
  \label{eq:13-proposal}
\end{equation}
with $\bm{\xi}$ a standard normal variate.  In atomic units $D=\tfrac12$,
which comes directly from the $\tfrac12$ in the kinetic energy operator; the
time step $\Delta t$ is a parameter to be tuned, exactly as $\Delta$ was.

\paragraph{Correcting for the asymmetry.}
The proposal is now \emph{not} symmetric: drifting towards high density makes
$\bm{x}\to\bm{y}$ and $\bm{y}\to\bm{x}$ different.  Metropolis is no longer
valid and must be replaced by Metropolis-Hastings, which was Exercise~3(d) of
Chapter~\ref{chap:statelements}.  The transition density is the Green's
function of Eq.~(\ref{eq:13-fokkerplanck}),
\begin{equation}
  G(\bm{y},\bm{x},\Delta t)
   = \frac{1}{(4\pi D\Delta t)^{3N/2}}
     \exp\left(-\frac{\left(\bm{y}-\bm{x}
       -D\Delta t\bm{F}(\bm{x})\right)^2}{4D\Delta t}\right),
  \label{eq:13-greensfunction}
\end{equation}
and the acceptance ratio becomes
\begin{equation}
  \boxed{\;
  q(\bm{y},\bm{x})
   = \frac{G(\bm{x},\bm{y},\Delta t)\,\left|\Psi_T(\bm{y})\right|^2}
          {G(\bm{y},\bm{x},\Delta t)\,\left|\Psi_T(\bm{x})\right|^2} \;}
  \label{eq:13-mhratio}
\end{equation}
The prefactors of $G$ cancel in the ratio, and expanding the exponents leaves
a form that needs no exponentials at all,
\begin{equation}
  \ln\frac{G(\bm{x},\bm{y})}{G(\bm{y},\bm{x})}
   = \frac{1}{2}\left(\bm{F}(\bm{x})+\bm{F}(\bm{y})\right)\cdot
     \left[\frac{D\Delta t}{2}
       \left(\bm{F}(\bm{x})-\bm{F}(\bm{y})\right)
       - \bm{y} + \bm{x}\right],
  \label{eq:13-greensratio}
\end{equation}
which is what the program evaluates.

\paragraph{The quantum force for our trial function.}
Differentiating Eq.~(\ref{eq:13-trial}),
\begin{equation}
  \bm{F}_1 = -2\alpha\omega\bm{r}_1
    + \frac{2\left(\bm{r}_1-\bm{r}_2\right)d^2}{r_{12}},
  \qquad
  \bm{F}_2 = -2\alpha\omega\bm{r}_2
    - \frac{2\left(\bm{r}_1-\bm{r}_2\right)d^2}{r_{12}},
  \label{eq:13-quantumforcedot}
\end{equation}
again with $d=1/(1+\beta r_{12})$.  It is a \emph{sum} of two terms: the first
pulls the electron towards the centre of the trap, the second pushes the pair
apart.

\begin{notebox}
\textbf{A wrong quantum force is not a wrong answer.} This deserves emphasis,
because it is a trap.  If Eq.~(\ref{eq:13-quantumforcedot}) is coded
incorrectly -- and the two terms are easy to combine wrongly -- the
calculation still converges to the right energy, provided the \emph{same}
expression is used both to propose the move and to evaluate the Green's
function.  Detailed balance holds for any drift; a bad drift merely proposes
bad moves.  Replacing the sum in Eq.~(\ref{eq:13-quantumforcedot}) by a
product, \texttt{vmc.py} still gets $3.000256\pm0.000399$, but the acceptance
rate falls from $87\%$ to $57\%$ and the correlation time doubles from $1.21$
to $2.48$.  The error is invisible in the answer and shows up only in the
efficiency, which is exactly why the quantum force is checked against finite
differences before use.  A wrong \emph{local energy}, by contrast, gives a
wrong number.
\end{notebox}

%----------------------------------------------------------------------
\section{Results}
\label{sec:vmcresults}

Table~\ref{tab:13-importance} scans the time step.  The comparison with
Table~\ref{tab:13-brute} is the point of the whole chapter.  At
$\Delta t=0.5$ the acceptance rate is $88\%$ while the walker is taking steps
that a uniform proposal of comparable size would have rejected almost always,
and the correlation time is $1.24$ against a best of $2.24$ for brute force.
There is no bias to worry about at large $\Delta t$: the Metropolis-Hastings
ratio corrects the discretisation of the Langevin equation exactly, so a
coarse time step costs accuracy only through the acceptance rate.

\begin{table}[h]
\centering
\renewcommand{\arraystretch}{1.25}
\setlength{\tabcolsep}{11pt}
\begin{tabular}{rrlll}
\toprule
$\Delta t$ & acceptance & energy & error & $\tau$\\
\midrule
$0.010$ & $100.0\%$ & $3.001377$ & $0.001969$ & $52.53$\\
$0.050$ & $\phantom{0}99.5\%$ & $3.001285$ & $0.000768$ & $\phantom{0}8.79$\\
$0.100$ & $\phantom{0}98.7\%$ & $3.000852$ & $0.000551$ & $\phantom{0}4.54$\\
$0.300$ & $\phantom{0}93.8\%$ & $3.000394$ & $0.000320$ & $\phantom{0}1.57$\\
$0.500$ & $\phantom{0}87.6\%$ & $3.000526$ & $0.000280$ & $\phantom{0}1.24$\\
$1.000$ & $\phantom{0}69.6\%$ & $3.000962$ & $0.000267$ & $\phantom{0}1.17$\\
\bottomrule
\end{tabular}
\caption{Importance-sampled variational Monte Carlo, same system and same
number of cycles as Table~\ref{tab:13-brute}.  Computed by \texttt{vmc.py}.}
\label{tab:13-importance}
\end{table}

Run at their respective best settings on equal numbers of cycles, the two
samplers give
\begin{center}
\renewcommand{\arraystretch}{1.2}
\setlength{\tabcolsep}{12pt}
\begin{tabular}{lrlrr}
\toprule
sampler & acceptance & energy $\pm$ error & $\tau$ & $n_{\mathrm{eff}}$\\
\midrule
brute force, $\Delta=4$ & $35.4\%$ & $3.000865\pm0.000265$ & $2.33$ & $25\,788$\\
importance, $\Delta t=0.5$ & $87.5\%$ & $3.000463\pm0.000199$ & $1.22$ & $49\,017$\\
\bottomrule
\end{tabular}
\end{center}
a factor $1.34$ in the error, which by the square-root law is a factor of
about $1.8$ in computer time.  That is a real but unspectacular gain, and it
is worth being honest about the scale: for a system this small, with a very
good trial function, brute force is already close to optimal.  The advantage
of importance sampling grows with the number of particles and with the
sharpness of the wave function, and it becomes indispensable in diffusion
Monte Carlo, where the drift is not an efficiency device but part of the
algorithm.

\paragraph{The variational surface.}
Table~\ref{tab:13-surface} scans both parameters.  The energy is remarkably
flat near the minimum -- the whole central block agrees to within a few parts
in $10^4$ -- while the variance varies by a factor of ten across the same
region.  This is the practical argument for variance minimisation: the
variance has a sharper minimum and a known target value, so it locates the
optimum more precisely than the energy does.

\begin{table}[h]
\centering
\renewcommand{\arraystretch}{1.2}
\setlength{\tabcolsep}{8pt}
\small
\begin{tabular}{rrrrrr}
\toprule
& \multicolumn{5}{c}{$\beta$}\\
\cmidrule(l){2-6}
$\alpha$ & $0.30$ & $0.36$ & $0.40$ & $0.44$ & $0.50$\\
\midrule
\multicolumn{6}{l}{\emph{energy}}\\
$0.94$ & $3.01804$ & $3.00780$ & $3.00460$ & $3.00316$ & $3.00357$\\
$0.96$ & $3.01232$ & $3.00415$ & $3.00197$ & $3.00127$ & $3.00266$\\
$0.98$ & $3.00786$ & $3.00142$ & $3.00043$ & $3.00081$ & $3.00295$\\
$1.00$ & $3.00492$ & $3.00033$ & $3.00012$ & $3.00145$ & $3.00445$\\
$1.02$ & $3.00294$ & $3.00025$ & $3.00096$ & $3.00297$ & $3.00671$\\
\midrule
\multicolumn{6}{l}{\emph{variance}}\\
$0.94$ & $0.03241$ & $0.01420$ & $0.00777$ & $0.00478$ & $0.00545$\\
$0.96$ & $0.02397$ & $0.00855$ & $0.00378$ & $0.00227$ & $0.00490$\\
$0.98$ & $0.01780$ & $0.00521$ & $0.00197$ & $0.00192$ & $0.00640$\\
$1.00$ & $0.01360$ & $0.00387$ & $0.00222$ & $0.00356$ & $0.00985$\\
$1.02$ & $0.01154$ & $0.00448$ & $0.00434$ & $0.00704$ & $0.01508$\\
\bottomrule
\end{tabular}
\caption{The variational surface, importance sampling with $\Delta t=0.5$ and
$15\,000$ cycles per point.  The energy is flat near its minimum; the variance
is not.  Computed by \texttt{vmc.py}.}
\label{tab:13-surface}
\end{table}

\paragraph{The final number.}
Two hundred thousand cycles at $\alpha=1.00$, $\beta=0.38$ give
\begin{equation}
  E_{\mathrm{VMC}} = 3.000603 \pm 0.000120,
  \qquad \sigma^2 = 0.00258,
  \qquad \tau = 1.12 ,
  \label{eq:13-final}
\end{equation}
against the exact $3$.  The gap of $6\times10^{-4}$ is five standard errors
wide, and that is the point rather than a problem.  A variational energy is an
upper bound, so what is being resolved here is not a statistical discrepancy
but the deficiency of a two-parameter ansatz -- and it is resolved only
because the statistical error has been pushed below it.  Removing the residual
requires a better trial function, or diffusion Monte Carlo, which projects it
out.

Set against Chapter~\ref{chap:applications}, where the best coupled-cluster
calculation in a $42$-orbital basis gave $3.013626$, the variational Monte
Carlo error is smaller by a factor of about twenty-three.  The reason is not
that the many-body treatment is better -- for two electrons CCSD is exact
within its basis -- but that there is no basis here at all.  The Jastrow
factor depends on $r_{12}$ directly and satisfies the cusp condition exactly,
and no finite sum of oscillator products can do either.

%----------------------------------------------------------------------
\section{Summary and what comes next}
\label{sec:vmcsummary}

A variational Monte Carlo program has four parts, and each has been dealt with
in turn: a trial wave function whose form is dictated by the physics
(Section~\ref{sec:vmctrial}); a local energy, differentiated in closed form
and verified (Section~\ref{sec:vmclocalenergy}); a sampler, which is the
Metropolis algorithm applied to $|\Psi_T|^2$
(Section~\ref{sec:vmcbrute}); and an error estimate that accounts for the
correlations a Markov chain introduces
(Section~\ref{sec:correlations}).  Importance sampling
(Section~\ref{sec:vmcimportance}) improves the third of these by proposing
moves along the quantum force and correcting the asymmetry with
Metropolis-Hastings.

Three things are still missing, and they are the subject of what follows.  The
error analysis here used the autocorrelation time; for production runs the
blocking method of Section~\ref{sec:correlations} is what one actually uses.
The parameters were located by scanning a grid, which does not scale beyond
two or three of them; gradient-based optimisation using the derivative of the
energy with respect to the parameters is the standard replacement.  And the
result remains an upper bound set by the quality of $\Psi_T$.  Diffusion Monte
Carlo removes that last limitation by letting the trial function evolve in
imaginary time towards the true ground state, using the same Fokker-Planck
machinery of Section~\ref{sec:vmcimportance} -- which is why it was worth
setting up carefully here.

%=============================================================================
\section{Exercises}
%=============================================================================

\subsection*{Exercise 1: the local energy}
\label{ex:13-localenergy}
\addcontentsline{toc}{subsection}{Exercise 1: the local energy}

\begin{enumerate}
  \item[a)] Derive the first line of Eq.~(\ref{eq:13-localenergydot}), the
  local energy of the pure Gaussian trial function.
  \item[b)] Show that $E_L$ is constant when $\alpha=1$ and the interaction is
  switched off, and say what its value is.
  \item[c)] Show that the $1/r_{12}$ terms in
  Eq.~(\ref{eq:13-localenergydot}) cancel as $r_{12}\to0$.
\end{enumerate}

\paragraph{Answer to a).}
With $\Psi=\exp(-\alpha\omega R^2/2)$ and $R^2=r_1^2+r_2^2$, use
$\nabla^2\Psi/\Psi = \nabla^2\ln\Psi + (\nabla\ln\Psi)^2$.  In two dimensions
$\nabla_i^2(-\alpha\omega r_i^2/2)=-2\alpha\omega$ for each particle, so the
first term contributes $-4\alpha\omega$, and
$(\nabla\ln\Psi)^2=\alpha^2\omega^2R^2$.  Hence
\[
  -\tfrac12\frac{\nabla^2\Psi}{\Psi}
   = 2\alpha\omega - \tfrac12\alpha^2\omega^2R^2 ,
\]
and adding the potential $\tfrac12\omega^2R^2 + 1/r_{12}$ gives
$2\alpha\omega + \tfrac12\omega^2(1-\alpha^2)R^2 + 1/r_{12}$.

\paragraph{Answer to b).}
Setting $\alpha=1$ kills the $R^2$ term and dropping $1/r_{12}$ leaves
$E_L=2\hbar\omega$, independent of position.  This is the zero-variance
principle in its simplest form: without the interaction the trial function is
the exact ground state of $\hat{H}_0$, its energy is $2\hbar\omega$, and a
single sample would suffice.  All the variance in the interacting problem
comes from the Coulomb term and from the imperfection of the Jastrow factor.

\paragraph{Answer to c).}
As $r_{12}\to0$, $d=1/(1+\beta r_{12})\to1$, so the term $-d^2/r_{12}$ tends
to $-1/r_{12}$ and cancels the $+1/r_{12}$ on the first line.  What survives
is finite: expanding to first order, $d^2\approx1-2\beta r_{12}$, and
\[
  \frac{1}{r_{12}}\left(1-d^2\right) \longrightarrow 2\beta ,
\]
so the singular parts cancel and leave a constant.  This is
Table~\ref{tab:13-cusp} in algebra, and it is why the numerator of the Jastrow
exponent had to be exactly $1$ and not a free parameter.

\subsection*{Exercise 2: the quantum force and the Green's function}
\label{ex:13-quantumforce}
\addcontentsline{toc}{subsection}{Exercise 2: the quantum force}

\begin{enumerate}
  \item[a)] Derive Eq.~(\ref{eq:13-quantumforcedot}) from
  Eq.~(\ref{eq:13-trial}).
  \item[b)] Derive the simplified Green's function ratio
  Eq.~(\ref{eq:13-greensratio}) from Eq.~(\ref{eq:13-greensfunction}).
  \item[c)] Show that Metropolis-Hastings with the ratio
  Eq.~(\ref{eq:13-mhratio}) satisfies detailed balance for \emph{any} drift
  $\bm{F}$, and explain what goes wrong if $\bm{F}$ is wrong.
\end{enumerate}

\paragraph{Answer to a).}
$\ln\Psi_T = -\tfrac12\alpha\omega(r_1^2+r_2^2) + r_{12}/(1+\beta r_{12})$.
The first term gives $\nabla_1 = -\alpha\omega\bm{r}_1$.  For the second, note
$\nabla_1 r_{12} = (\bm{r}_1-\bm{r}_2)/r_{12}$ and
\[
  \frac{d}{dr_{12}}\frac{r_{12}}{1+\beta r_{12}}
   = \frac{1}{(1+\beta r_{12})^2} = d^2 ,
\]
so $\nabla_1[\,\cdot\,] = (\bm{r}_1-\bm{r}_2)d^2/r_{12}$.  Multiplying the sum
by two gives Eq.~(\ref{eq:13-quantumforcedot}); $\bm{F}_2$ follows from
$\nabla_2 r_{12} = -(\bm{r}_1-\bm{r}_2)/r_{12}$.

\paragraph{Answer to b).}
Write $\bm{u}=\bm{y}-\bm{x}$, $\bm{a}=D\Delta t\bm{F}(\bm{y})$ and
$\bm{b}=D\Delta t\bm{F}(\bm{x})$.  Then
\[
  \ln\frac{G(\bm{x},\bm{y})}{G(\bm{y},\bm{x})}
   = -\frac{(\bm{u}+\bm{a})^2 - (\bm{u}-\bm{b})^2}{4D\Delta t}
   = -\frac{2\bm{u}\cdot(\bm{a}+\bm{b}) + a^2 - b^2}{4D\Delta t}.
\]
Substituting back and factorising $a^2-b^2=(\bm{a}+\bm{b})\cdot
(\bm{a}-\bm{b})$ gives Eq.~(\ref{eq:13-greensratio}).  The prefactors of $G$
are identical in numerator and denominator and cancel, which is why the
normalisation of the Green's function never has to be computed.

\paragraph{Answer to c).}
Detailed balance requires
$T(\bm{x}\to\bm{y})A(\bm{x}\to\bm{y})P(\bm{x})
= T(\bm{y}\to\bm{x})A(\bm{y}\to\bm{x})P(\bm{y})$.  With $T=G$ and
$A=\min(1,q)$ from Eq.~(\ref{eq:13-mhratio}) this holds identically,
whatever $\bm{F}$ is, because $q$ is by construction the ratio that makes it
hold.  A wrong $\bm{F}$ therefore leaves the stationary distribution
untouched and the energy correct.  What it damages is the acceptance rate:
the drift proposes moves the wave function does not like, more of them are
rejected, and the correlation time rises.  The consistency requirement is that
the same $\bm{F}$ be used in the proposal and in $G$; using different ones
would break detailed balance and give a wrong answer.

\subsection*{Exercise 3: the cusp condition}
\label{ex:13-cusp}
\addcontentsline{toc}{subsection}{Exercise 3: the cusp condition}

\begin{enumerate}
  \item[a)] Derive Eq.~(\ref{eq:13-cuspderivation}) by keeping only the terms
  in the relative-motion local energy that survive as $r_{12}\to0$.
  \item[b)] Why is the coefficient different for parallel and antiparallel
  spins?
  \item[c)] Show that the Pad\'e-Jastrow form~(\ref{eq:13-padejastrow})
  reproduces the cusp at short range and saturates at long range, and say what
  goes wrong if the bare exponential is used instead.
\end{enumerate}

\paragraph{Answer to a).}
In the relative coordinate the kinetic operator contains
$d^2/dr^2 + (2/r)\,d/dr$ in two dimensions, and the potential contains
$1/r$.  For $E_L$ to stay finite as $r\to0$ the $1/r$ terms must cancel:
\[
  \frac{2}{r}\frac{{\cal R}'}{\cal R} + \frac{2}{r} = 0
  \quad\Longrightarrow\quad
  \frac{{\cal R}'}{\cal R}\bigg|_{r\to0} = \text{constant},
\]
which is Eq.~(\ref{eq:13-cuspderivation}).  Everything else -- the harmonic
term, the energy, the centre-of-mass part -- is regular at the origin and
cannot participate.

\paragraph{Answer to b).}
For parallel spins the spatial wave function is antisymmetric and already
vanishes linearly as $r_{12}\to0$, so the Coulomb singularity is partly
screened by the Pauli principle and less is asked of the correlation factor.
The angular momentum in the relative motion is $l=1$ rather than $l=0$, and
the factor $1/2(l+1)$ in Eq.~(\ref{eq:13-cuspderivation}) halves the
coefficient.

\paragraph{Answer to c).}
Expanding $ar/(1+\beta r) = ar(1-\beta r + \cdots)$ shows that the exponent is
$ar$ to first order, so the logarithmic derivative at the origin is $a$ and
the cusp condition is satisfied for any $\beta$.  As $r\to\infty$ the exponent
tends to $a/\beta$, a constant, so the factor saturates and the wave function
retains the decay of the mean-field part.  With the bare $\exp(ar)$ the
correlation factor grows without bound and overwhelms the Gaussian at large
separation: the trial function is not normalisable in the relevant sense and
the electrons are driven apart.  This is exactly what
Chapter~\ref{chap:statelements} warned about -- setting $\beta=0$ does not
switch the Jastrow factor off, it leaves $\exp(r_{12})$.

\subsection*{Exercise 4: numerical explorations}
\addcontentsline{toc}{subsection}{Exercise 4: numerical explorations}

\begin{enumerate}
  \item[a)] Using \texttt{vmc.py}, reproduce Tables~\ref{tab:13-brute}
  and~\ref{tab:13-importance}.  Plot the correlation time against the
  acceptance rate for both samplers on the same axes.  Is there a universal
  optimal acceptance rate?
  \item[b)] Replace the analytic local energy by the finite-difference one and
  compare the cost and the statistical error.  How small may the step $h$ be
  made before round-off dominates?
  \item[c)] Implement blocking, as in Section~\ref{sec:correlations}, and
  compare the error estimate with the one from the correlation time for a run
  at $\Delta t=0.01$, where $\tau$ is large.
  \item[d)] Repeat the calculation at $\hbar\omega=0.5$ and $\hbar\omega=0.1$.
  Where is the correlation strongest, and how do the optimal $\alpha$ and
  $\beta$ move?  Compare with Table~\ref{tab:11-omega}.
  \item[e)] Compute the derivative of the energy with respect to a variational
  parameter $\theta$ using
  \[
    \frac{\partial E}{\partial\theta}
     = 2\left(\mean{E_L\,\partial_\theta\ln\Psi_T}
       - \mean{E_L}\mean{\partial_\theta\ln\Psi_T}\right),
  \]
  and use it to optimise $\alpha$ and $\beta$ by gradient descent.  How many
  iterations does it take, and how does the cost compare with the grid scan of
  Table~\ref{tab:13-surface}?
  \item[f)] Extend the code to six electrons using a Slater determinant of the
  lowest three oscillator orbitals times the Pad\'e-Jastrow factor, and
  compare with the coupled-cluster result of
  Table~\ref{tab:11-hierarchy}.
\end{enumerate}
